\documentclass[11pt]{article}

\usepackage[margin=1in]{geometry}
\usepackage{amsmath,amssymb}
\usepackage{natbib}
\usepackage{graphicx}
\usepackage{booktabs}
\usepackage{xcolor}
\usepackage{hyperref}
\usepackage{cleveref}

\newcommand{\GMSL}{\text{GMSL}}
\newcommand{\GMST}{\text{GMST}}
\newcommand{\daT}{\mathrm{d}\alpha/\mathrm{d}T}

\title{Observationally Calibrated Sea Level Rise Projections to 2150}

\author{B.\ Minchew}
\date{Working Draft --- \today}

\begin{document}
\maketitle

% ====================================================================
% TARGET: Nature Communications, Nature Geoscience, or GRL
% LENGTH: ~4000 words (NatComms) or ~3500 (GRL with supplement)
% FIGURES: 4-5 main, supplementary as needed
% TIMELINE: Submittable within 4-6 weeks (analysis exists)
% ====================================================================

\begin{abstract}
% ~150 words
We present global mean sea level (GMSL) projections to 2150 derived from an observationally calibrated rate--temperature relationship.
Dynamic ordinary least squares (DOLS) applied to seven GMSL records and four global mean surface temperature (GMST) products yields a robust quadratic relationship between the rate of sea level rise and warming: the thermodynamic transient sensitivity $\daT = 2.85 \pm 0.38$~mm/yr/${}^\circ$C$^2$ across the thermodynamic dataset ensemble.
Projections under SSP1-2.6 through SSP5-8.5, produced by evaluating this relationship at FaIR-emulated temperature trajectories, exceed IPCC AR6 process-model-based thermodynamic projections by [X]\% under SSP2-4.5 and [Y]\% under SSP5-8.5 at 2100.
The difference arises because CMIP-class ocean models produce a genuinely more linear rate--temperature relationship, with transient sensitivity approximately half the observational value.
The observationally calibrated projection decomposes naturally into a thermodynamic component (constrained by the DOLS calibration) and an ice-sheet residual, providing an independent bound on cryospheric contributions.
\end{abstract}


% ====================================================================
\section{Introduction}
% ~500 words
% ====================================================================

% Frame as: we produce projections, not critique existing ones.
% 
% Key points:
% 1. SLR projections serve decisions from infrastructure (20 yr) to policy (100+ yr).
%    Decision-makers need projections anchored to observations.
%
% 2. Current projections (IPCC AR6, FACTS) rely on process models that lack
%    systematic observational calibration at the aggregate level. They calibrate
%    component models individually but do not calibrate the total against the
%    observed rate-temperature relationship.
%
% 3. Semi-empirical models (Rahmstorf 2007, Vermeer & Rahmstorf 2009) calibrate
%    against observations but use simple functional forms and do not decompose
%    the signal into physically distinct components.
%
% 4. We present a method that combines the strengths of both: observational
%    calibration (like semi-empirical) with component decomposition (like FACTS).
%    The DOLS rate-temperature relationship provides the thermodynamic constraint;
%    the residual isolates the ice-sheet contribution.
%
% 5. We produce projection distributions under all SSPs to 2150 with uncertainty
%    that includes dataset structural uncertainty, and compare to IPCC AR6.


% ====================================================================
\section{Observational Calibration}
% ~800 words
% ====================================================================

\subsection{Dynamic OLS rate--temperature model}
% 
% Present the DOLS model:
%   rate(t) = alpha_0 * T(t) + (d_alpha/dT) * T(t)^2 + [optional SAOD term]
%
% Explain: DOLS adds leads/lags of dT/dt to absorb short-run dynamics,
% then discards them. HAC (Newey-West) standard errors for serial correlation.
% WLS with weights = 1/sigma_rate^2.
%
% This is a methods paragraph, not a results paragraph. Keep it concise.

\subsection{Multi-dataset robustness}
%
% The 6x4 robustness matrix: 6 GMSL records x 4 GMST products = 24 fits.
% Present Table 1 (the robustness matrix).
%
% Key result: d_alpha/dT > 0 robustly across all datasets except Dangendorf
% sterodynamic (which excludes the cryospheric signal by construction).
%
% Thermodynamic ensemble (8 pairs): d_alpha/dT = 2.85 +/- 0.38
% All-dataset ensemble (24 pairs): d_alpha/dT = 1.83 +/- 1.09
%
% GMSL choice dominates spread; GMST choice has modest impact.
% This quantifies the structural uncertainty from dataset choice.

\subsection{Epoch dependence}
%
% Brief mention of the start-date sensitivity (Frederikse, Dangendorf 
% at native vs 1950 start). alpha_0 and d_alpha/dT trade off.
% The sliding-window DOLS shows smooth evolution.
% This motivates using the 1950-start calibration as the primary estimate
% (most datasets available, stable estimates) and reporting native-start
% as a sensitivity check.
%
% Full sliding-window analysis in supplement.


% ====================================================================
\section{Projections}
% ~1000 words
% ====================================================================

\subsection{Method}
%
% Given DOLS coefficients (alpha_0, d_alpha/dT) and their covariance matrix,
% and given FaIR temperature trajectories T(t) for each SSP:
%
% 1. Draw (alpha_0, d_alpha/dT) from bivariate normal posterior.
% 2. Compute rate(t) = alpha_0 * T(t) + d_alpha/dT * T(t)^2.
% 3. Integrate to get cumulative GMSL: H(t) = integral_t0^t rate(tau) d_tau.
% 4. Add structural uncertainty from multi-dataset spread (outer loop).
%
% Produce N_MC = 10,000 trajectories per SSP.
% Report median and 17-83% (likely) and 5-95% (very likely) ranges.

\subsection{DOLS projections}
%
% Present Table 2: DOLS projection ranges at 2050, 2100, 2150 for each SSP.
% Format to match IPCC/FACTS Table A1 for direct comparison.
%
% Present Figure 1: Projection fan chart (time vs GMSL) for all SSPs.
% Median lines + 17-83% shading. 
% Overlay observed GMSL (satellite altimetry) for hindcast validation.

\subsection{Decomposition into thermodynamic and residual}
%
% The DOLS model captures the "thermodynamic" component: everything that
% responds to GMST through the rate-temperature relationship. This includes
% thermal expansion, glaciers, Greenland SMB — the components that scale
% with temperature.
%
% The ice-sheet dynamic component (WAIS, EAIS discharge, GrIS discharge)
% is NOT captured by the DOLS model because it responds to ocean thermal
% forcing, not GMST directly.
%
% Define: GMSL_residual(t) = GMSL_obs(t) - GMSL_DOLS(t)
% During the calibration period, this residual represents the non-thermodynamic
% contribution. Its magnitude and trend provide an observationally-derived
% bound on the ice-sheet dynamic contribution.
%
% Figure 2: Residual time series, 1950-2020. Compare to IMBIE AIS+GrIS_discharge.
% The residual should be small and consistent with known ice-sheet contributions.


% ====================================================================
\section{Comparison with IPCC AR6 / FACTS}
% ~600 words
% ====================================================================

%
% Figure 3: Side-by-side comparison of DOLS projections vs FACTS projections.
% Two panels or overlaid:
%   - DOLS median + 17-83% range
%   - FACTS Workflow 2e median + 17-83% range
% For each SSP.
%
% Key findings:
% 1. At 2100 under SSP2-4.5: DOLS projects [X] mm vs FACTS [Y] mm.
%    The DOLS thermodynamic projection is higher because the observational
%    quadratic sensitivity exceeds the IPCC linear sensitivity.
%
% 2. The discrepancy grows with warming: small under SSP1-2.6, large under SSP5-8.5.
%    This is because the quadratic term contributes more at higher T.
%
% 3. The IPCC process-model thermodynamic sensitivity is genuinely more linear.
%    Present the comparison: fit rate-T to IPCC thermodynamic projections,
%    find alpha_0 ~ 1.9-2.5, no significant quadratic. Power analysis confirms
%    this is not a sample-size artifact. Observational quadratic falls outside
%    IPCC 95% CI for SSP2-4.5 through SSP5-8.5.
%
% 4. Physical mechanism: Kuhlbrodt & Gregory (2012) — weak ocean stratification
%    in CMIP models absorbs heat too efficiently into deep ocean.
%
% Frame this comparison as: "the DOLS projections provide an independent
% observational check on process-model projections. Where they disagree,
% the disagreement is informative about structural model deficiencies."


% ====================================================================
\section{Discussion}
% ~500 words
% ====================================================================

\subsection{Strengths and limitations}
%
% Strengths:
% - Anchored to observations, not process-model structure
% - Robustly calibrated across datasets
% - Transparent uncertainty (coefficient + structural)
% - Directly produces projection distributions in IPCC-comparable format
%
% Limitations:
% - Assumes rate-temperature relationship is stationary. Sliding-window
%   analysis suggests it may be time-varying. This is a known limitation
%   that could be addressed with a time-varying coefficient model.
% - Cannot resolve component-specific dynamics (AIS response to ocean forcing,
%   GrIS discharge, MISI). The residual captures these but cannot project them.
% - The quadratic form is empirical. At very high T (>4°C above pre-industrial),
%   saturation effects (glacier volume depletion, thermosteric equilibrium)
%   may cause the relationship to deviate from quadratic.

\subsection{Implications for AR7}
%
% The DOLS projections are directly usable by AR7 as an independent line of
% evidence alongside process-model projections. They provide:
% 1. A check on whether process-model thermodynamic projections are consistent
%    with the observed rate-temperature relationship.
% 2. An observationally-derived lower bound on total GMSL (the DOLS projection
%    excludes the ice-sheet dynamic contribution, so it is a lower bound on
%    the total if ice-sheet dynamics contribute positively).
% 3. Projection distributions in the same format as FACTS for direct comparison.


% ====================================================================
\section{Data and Code Availability}
% ====================================================================
% Projection distributions (CSV): Zenodo deposit
% Analysis code: GitHub repository
% All input data: publicly available (list DOIs)


% ====================================================================
% SUPPLEMENTARY MATERIALS
% ====================================================================
% SM1: Full DOLS methodology (lag selection, kernel specification, HAC)
% SM2: Multi-dataset robustness matrix (full table with all 24 fits)
% SM3: Start-date sensitivity analysis
% SM4: Sliding-window DOLS coefficient evolution
% SM5: SAOD sensitivity
% SM6: IPCC emergent sensitivity comparison (all 5 SSPs, power analysis)
% SM7: FaIR temperature trajectory specification
% SM8: Full projection tables at decadal intervals

% FIGURES:
% Main Fig 1: Projection fan chart (all SSPs, with observed GMSL overlay)
% Main Fig 2: DOLS residual time series vs IMBIE ice-sheet contributions
% Main Fig 3: DOLS vs FACTS comparison
% Main Fig 4: Rate-temperature relationship: observational vs IPCC
% Fig S1: DOLS robustness matrix visualization
% Fig S2: Start-date sensitivity
% Fig S3: Sliding-window coefficient evolution
% Fig S4: IPCC sensitivity per SSP with power analysis
% Fig S5: Projection distributions at 2100 (density plots per SSP)

% TABLES:
% Table 1: Multi-dataset DOLS robustness matrix (main text, condensed)
% Table 2: Projection ranges at 2050, 2100, 2150 (main text)
% Table S1: Full robustness matrix (all 24 fits)
% Table S2: Start-date sensitivity
% Table S3: IPCC sensitivity comparison (all 5 SSPs)
% Table S4: Full projection tables (decadal, all SSPs)


\bibliographystyle{agsm}
\bibliography{slr_references}

\end{document}
